\begin{threeparttable}
\begin{tabular}{lccc}
\toprule
& \multicolumn{3}{c}{\textbf{Dep. var.}: Five-year log change in} \\
\cmidrule(lr){2-4}
& Borrowing firms & Loans per borrower & Credit per loan \\
\midrule
$\widehat{\Delta_5\ln(\text{C})}$
& 0.161\sym{***} & 0.202\sym{***} & 0.637\sym{***} \\
& (0.035) & (0.034) & (0.044) \\
& [0.197] & [$<$0.001] & [$<$0.001] \\
\midrule
FS Coefficient & 1.104 & 1.104 & 1.104 \\
               & (0.063) & (0.063) & (0.063) \\
FS F & 310.25 & 310.25 & 310.25 \\
Mean & 670 & 1.95 & 455 \\
Obs. & 1{,}390 & 1{,}390 & 1{,}390 \\
Markets & 556 & 556 & 556 \\
\bottomrule
\end{tabular}
\begin{tablenotes}[flushleft]
\footnotesize
\item \textit{Notes:} Market credit satisfies $C = \text{borrowing firms}
\times \text{loans per borrower} \times \text{credit per loan}$, so the
three coefficients sum to the response of credit to itself, one, by
construction ($0.161 + 0.202 + 0.637 = 1.00$). Outcomes are from the credit
registry.  Means in levels: the average market-year in the estimation sample has
$670$ borrowing firms with $1.95$ loans per borrower. To scale the credit
response itself we use the median market, whose credit stock is $4.5$
million USD: a ten percent credit supply shock corresponds to roughly $455$
thousand USD of additional credit in the median market. See Table~\ref{main_table} for specification and inference
details; wild cluster bootstrap $p$-values in brackets.
\sym{*}~$p<0.10$, \sym{**}~$p<0.05$, \sym{***}~$p<0.01$.

\end{tablenotes}
\end{threeparttable}